% GENERATED FILE. Edit canonical lessons, then run npm run generate:books.
% canonical-source-hash: fnv1a64:1bf24f4b65f94db9
% canonical-lessons: LA-C485-imber, LA-C485-madeo, LA-C485-pluvialis, LA-C485-nimbosus, LA-C485-turbidus

\chapter{Downpour, To be soaking wet, Rainy, Stormy, Turbulent}
\label{ch:la-downpour-to-be-soaking-wet-rainy-stormy-turbulent}
\hlchaptermodality{485}

\begin{chapteropening}
\textbf{By the end of this chapter:} \emph{I can describe rain and storms: a downpour, being soaked, rainy and stormy weather, and turbulent water.}

Complete the last lesson of chapter 485: I can describe rain and storms: a downpour, being soaked, rainy and stormy weather, and turbulent water.
\end{chapteropening}

\section[imber, imbris]{imber, imbris — a downpour, a heavy shower of rain}

\label{lesson:LA-C485-imber}



\begin{quote}
Before the new one: say the Latin for an ash tree, then the Latin for shady.
\end{quote}

\subsection*{You'll want to know: imber, imbris}
\textbf{imber, imbris} — \textquotedblleft{}a downpour, a heavy shower of rain\textquotedblright{}.

\begin{quote}
Subitō imber magnus cadit, et viātōrēs ad caupōnam currunt.
\end{quote}

\begin{quote}
\emph{Suddenly a heavy downpour falls, and the travellers run to the inn.}
\end{quote}

Third declension, masculine: \emph{imber, imbris m.}, genitive plural \emph{imbrium}. Heavier than \emph{pluvia}.

\begin{cousinweb}
\emph{imber} is inherited; Lewis \& Short list Sanskrit \emph{abhra} 'cloud', Greek \emph{ombros} 'rain' and Latin \emph{umbra} 'shade' as kin. EtymDB lists it without a parent.
\end{cousinweb}

\begin{grammarlens}[title={What you've built}]
One more word for telling, explaining and describing at B1.
\end{grammarlens}

\subsection*{Your turn}
\begin{itemize}
\raggedright
  \item \emph{Say it:} \emph{imber, imbris}
  \item \emph{Say it:} \emph{imber, imbris}, once more
  \item \emph{Say it:} a sentence of your own with \emph{imber, imbris}
  \item \emph{Recall:} say the Latin for an ash tree, then the Latin for shady, then say \emph{imber, imbris} again
\end{itemize}

\begin{tcolorbox}[breakable,colback=teal!4,colframe=teal!35!black,title={Before you move on}]
What does imber, imbris mean? (\textquotedblleft{}A downpour, a heavy shower of rain\textquotedblright{}.) Say it once more.
\end{tcolorbox}
\section[madeō, madēre]{madeō, madēre — to be soaking wet, to drip}

\label{lesson:LA-C485-madeo}



\begin{quote}
Before the new one: say the Latin for shady, then the Latin for a downpour.
\end{quote}

\subsection*{You'll want to know: madeō, madēre}
\textbf{madeō, madēre} — \textquotedblleft{}to be soaking wet, to drip\textquotedblright{}.

\begin{quote}
Post imbrem vestēs viātōrum madent.
\end{quote}

\begin{quote}
\emph{After the downpour the travellers' clothes are soaking wet.}
\end{quote}

\emph{madeō, madēre, maduī}, second conjugation, no supine. Also, in comedy, to be drunk.

\begin{cousinweb}
\emph{madeō} is inherited from PIE \emph{meh\textsubscript{2}d-} 'to drip, to ooze' (EtymDB); Lewis \& Short compare Greek \emph{madaō} 'to be moist'.
\end{cousinweb}

\begin{grammarlens}[title={What you've built}]
One more word for telling, explaining and describing at B1.
\end{grammarlens}

\subsection*{Your turn}
\begin{itemize}
\raggedright
  \item \emph{Say it:} \emph{madeō, madēre}
  \item \emph{Say it:} \emph{madeō, madēre}, once more
  \item \emph{Say it:} a sentence of your own with \emph{madeō, madēre}
  \item \emph{Recall:} say the Latin for shady, then the Latin for a downpour, then say \emph{madeō, madēre} again
\end{itemize}

\begin{tcolorbox}[breakable,colback=teal!4,colframe=teal!35!black,title={Before you move on}]
What does madeō, madēre mean? (\textquotedblleft{}To be soaking wet, to drip\textquotedblright{}.) Say it once more.
\end{tcolorbox}
\section[pluviālis, pluviāle]{pluviālis, pluviāle — rainy, of rain, bringing rain}

\label{lesson:LA-C485-pluvialis}



\begin{quote}
Before the new one: say the Latin for a downpour, then the Latin for to be soaking wet.
\end{quote}

\subsection*{You'll want to know: pluviālis, pluviāle}
\textbf{pluviālis, pluviāle} — \textquotedblleft{}rainy, of rain, bringing rain\textquotedblright{}.

\begin{quote}
Ventus pluviālis ā marī venit.
\end{quote}

\begin{quote}
\emph{The rainy wind comes from the sea.}
\end{quote}

An adjective of two endings, from \emph{pluvia} 'rain'. \emph{Aqua pluviālis}, rainwater. English \emph{pluvial} comes from it.

\begin{cousinweb}
\emph{pluviālis} is built on \emph{pluvia} 'rain' with the suffix \emph{-ālis} (EtymDB; Lewis \& Short), and \emph{pluvia} comes from \emph{pluit} 'it rains'.
\end{cousinweb}

\begin{grammarlens}[title={What you've built}]
One more word for telling, explaining and describing at B1.
\end{grammarlens}

\subsection*{Your turn}
\begin{itemize}
\raggedright
  \item \emph{Say it:} \emph{pluviālis, pluviāle}
  \item \emph{Say it:} \emph{pluviālis, pluviāle}, once more
  \item \emph{Say it:} a sentence of your own with \emph{pluviālis, pluviāle}
  \item \emph{Recall:} say the Latin for a downpour, then the Latin for to be soaking wet, then say \emph{pluviālis, pluviāle} again
\end{itemize}

\begin{tcolorbox}[breakable,colback=teal!4,colframe=teal!35!black,title={Before you move on}]
What does pluviālis, pluviāle mean? (\textquotedblleft{}Rainy, of rain, bringing rain\textquotedblright{}.) Say it once more.
\end{tcolorbox}
\section[nimbōsus, nimbōsa, nimbōsum]{nimbōsus, nimbōsa, nimbōsum — stormy, full of rain clouds}

\label{lesson:LA-C485-nimbosus}



\begin{quote}
Before the new one: say the Latin for to be soaking wet, then the Latin for rainy.
\end{quote}

\subsection*{You'll want to know: nimbōsus, nimbōsa, nimbōsum}
\textbf{nimbōsus, nimbōsa, nimbōsum} — \textquotedblleft{}stormy, full of rain clouds\textquotedblright{}.

\begin{quote}
Caelum nimbōsum est; nautae in portū manent.
\end{quote}

\begin{quote}
\emph{The sky is stormy; the sailors stay in the harbour.}
\end{quote}

An adjective of three endings, from \emph{nimbus} 'rain cloud' with \emph{-ōsus} 'full of'.

\begin{cousinweb}
\emph{nimbōsus} is built on \emph{nimbus} 'rain cloud, storm' with the suffix \emph{-ōsus} 'full of' (EtymDB; Lewis \& Short).
\end{cousinweb}

\begin{grammarlens}[title={What you've built}]
One more word for telling, explaining and describing at B1.
\end{grammarlens}

\subsection*{Your turn}
\begin{itemize}
\raggedright
  \item \emph{Say it:} \emph{nimbōsus, nimbōsa, nimbōsum}
  \item \emph{Say it:} \emph{nimbōsus, nimbōsa, nimbōsum}, once more
  \item \emph{Say it:} a sentence of your own with \emph{nimbōsus, nimbōsa, nimbōsum}
  \item \emph{Recall:} say the Latin for to be soaking wet, then the Latin for rainy, then say \emph{nimbōsus, nimbōsa, nimbōsum} again
\end{itemize}

\begin{tcolorbox}[breakable,colback=teal!4,colframe=teal!35!black,title={Before you move on}]
What does nimbōsus, nimbōsa, nimbōsum mean? (\textquotedblleft{}Stormy, full of rain clouds\textquotedblright{}.) Say it once more.
\end{tcolorbox}
\section[turbidus, turbida, turbidum]{turbidus, turbida, turbidum — turbulent, wild; cloudy, muddy (of water)}

\label{lesson:LA-C485-turbidus}



\begin{quote}
Before the new one: say the Latin for rainy, then the Latin for stormy.
\end{quote}

\subsection*{You'll want to know: turbidus, turbida, turbidum}
\textbf{turbidus, turbida, turbidum} — \textquotedblleft{}turbulent, wild; cloudy, muddy (of water)\textquotedblright{}.

\begin{quote}
Post tempestātem flūmen turbidum est.
\end{quote}

\begin{quote}
\emph{After the storm the river is cloudy and wild.}
\end{quote}

An adjective of three endings, from \emph{turba} 'disorder'. Of water, weather or a crowd. English \emph{turbid} comes from it.

\begin{cousinweb}
\emph{turbidus} is built on \emph{turba} 'disorder, commotion' with the suffix \emph{-idus} (EtymDB; Lewis \& Short). English \emph{turbid} comes from it.
\end{cousinweb}

\begin{grammarlens}[title={What you've built}]
One more word for telling, explaining and describing at B1.
\end{grammarlens}

\subsection*{Your turn}
\begin{itemize}
\raggedright
  \item \emph{Say it:} \emph{turbidus, turbida, turbidum}
  \item \emph{Say it:} \emph{turbidus, turbida, turbidum}, once more
  \item \emph{Say it:} a sentence of your own with \emph{turbidus, turbida, turbidum}
  \item \emph{Recall:} say the Latin for rainy, then the Latin for stormy, then say \emph{turbidus, turbida, turbidum} again
\end{itemize}

\begin{tcolorbox}[breakable,colback=teal!4,colframe=teal!35!black,title={Before you move on}]
What does turbidus, turbida, turbidum mean? (\textquotedblleft{}Turbulent, wild; cloudy, muddy (of water)\textquotedblright{}.) Say it once more.
\end{tcolorbox}
